\bookchapter{الفصل 39: هندسة السياق — اجعل الذكاء الاصطناعي يفهم نيتك}

\chapterintro{هل جربت هذا الموقف: تفتح نافذة محادثة جديدة، وكأن الذكاء الاصطناعي أصيب بفقدان الذاكرة ونسي كل شيء عن مشروعك؟ ليس لأنه يتعمد إزعاجك، بل لأنه «سريع النسيان» بطبيعته. يعلمك هذا الفصل كيف تركّب له «نظام ذاكرة طويلة المدى»، وتتعرف على أحد أشهر نماذج التطوير في 2026 — التطوير المدفوع بالمواصفات.}

\booksection{أولاً: ملف ذاكرة المشروع: «الذاكرة طويلة المدى» للذكاء الاصطناعي}

للنماذج الذكية قيد جوهري: ذاكرتها «قصيرة المدى». في بداية كل محادثة، لا يعرف الذكاء الاصطناعي شيئًا عن مشروعك. عليك إخباره باستمرار: «نستخدم React» و«الخادم الخلفي Node.js» و«قاعدة البيانات PostgreSQL» و«نمط الكود يستخدم علامتي اقتباس مزدوجتين». هذا غير فعال ويسهل فيه النسيان — كما لو كنت تعيد تدريب زميل جديد كل يوم.

الحل هو «ملف ذاكرة المشروع» — ملف نصي في جذر المشروع، تقرأه أدوات الذكاء الاصطناعي تلقائيًا عند بداية كل محادثة. لكل أداة تسميتها الخاصة:

\begin{center}
\begin{tabularx}{\linewidth}{|X|X|}
\hline
\textbf{الأداة} & \textbf{اسم الملف} \tabularnewline
\hline
Cursor & \texttt{.\allowbreak cursorrules} \tabularnewline
Claude Code & \texttt{CLAUDE.\allowbreak md} \tabularnewline
GitHub Copilot & \texttt{copilot-\allowbreak instructions.\allowbreak md} \tabularnewline
الحل العام & \texttt{AGENTS.\allowbreak md} \tabularnewline
\hline
\end{tabularx}
\end{center}

\begin{center}\small\color{gray}تسميات ملفات ذاكرة المشروع لكل أداة\end{center}

رغم اختلاف التسميات، المبدأ واحد: اكتب في هذا الملف المعلومات التي تريد أن «يتذكرها» الذكاء الاصطناعي دائمًا. ملف الذاكرة الجيد يشمل: نظرة عامة على المشروع (ما هو)، والمكدس التقني (ما التقنيات المستخدمة)، ومعايير الكود (التسمية والمسافة البادئة ونمط التعليقات)، وبنية المجلدات (كيف يُنظم الكود)، واتفاقيات التطوير (استراتيجية الفروع، وقواعد الإيداع، ومتطلبات الاختبار). ملف CLAUDE.md لـ Karpathy 70 سطرًا فقط، لكنه يغطي كل المعلومة المهمة — مختصر ومكتمل.

\begin{lstlisting}
# المشروع: تطبيق لتتبع المصروفات الشخصية
## نظرة عامة
تطبيق متكامل (full-stack) لتتبع المصروفات الشخصية يساعد المستخدم على تسجيل الإيرادات والمصروفات اليومية.
## الحزمة التقنية
- الواجهة الأمامية: React + TypeScript + Tailwind CSS
- الواجهة الخلفية: Node.js + Express + TypeScript
- قاعدة البيانات: SQLite (تطوير) / PostgreSQL (إنتاج)
- الاختبارات: Vitest
## قواعد الكود
- استخدم علامات الاقتباس المزدوجة
- مسافة بادئة بمسافتين
- أسماء الدوال بصيغة camelCase
- تعريفات الأنواع بصيغة PascalCase
- كل دالة يجب أن تملك تعليق JSDoc
## هيكل المجلدات
/src
  /frontend  - كود الواجهة الأمامية
  /backend   - كود الواجهة الخلفية
  /shared    - الأنواع المشتركة بين الواجهتين
  /tests     - ملفات الاختبار
## اتفاقيات التطوير
- رسائل الـ commit بالعربية
- كل ميزة جديدة تبدأ باختبار أولًا
- ردود الـ API بصيغة موحدة { success, data, error }
\end{lstlisting}

\begin{tcbwarn}{ملاحظة: لا تحوّله إلى «مستند متطلبات ثانٍ»}
انتبه لطول الملف — لا تكتبه كعشرات الصفحات من مستند المتطلبات. الهدف هو «تمكين الذكاء الاصطناعي من بناء فهم سريع للخلفية»، وليس «وصف كل وظيفة بالتفصيل». أبقِه مختصرًا ومنظّمًا، واكتب فقط «ما يحتاجه الذكاء الاصطناعي معرفته».

\end{tcbwarn}

\booksection{ثانيًا: التطوير المدفوع بالمواصفات (SDD)}

التطوير المدفوع بالمواصفات (Spec-Driven Development، SDD) هو نموذج تطوير نضج تدريجيًا في 2026، وفكرته في جملة: اكتب المواصفات أولًا، ثم الكود. بخلاف TDD التقليدي (التطوير المدفوع بالاختبارات)، لا تقتصر «المواصفات» في SDD على حالات الاختبار، بل تشمل وصف المتطلبات وتعريفات الواجهات ونماذج البيانات. يولد الذكاء الاصطناعي الكود بناءً على هذه المواصفات، ودور الإنسان هو كتابتها وصيانتها.

سير عمل SDD في خمس خطوات:

\begin{enumerate}
  \item صِف المتطلبات باللغة الطبيعية (في ملف .md)؛
  \item حدد الواجهات وهياكل البيانات (في ملف المواصفات)؛
  \item دع الذكاء الاصطناعي يولّد الكود وفقًا للمواصفات؛
  \item شغّل الاختبارات للتحقق من مطابقة المخرجات للمواصفات؛
  \item إذا لم تتطابق، عدّل المواصفات أو اطلب من الذكاء الاصطناعي تصحيح الكود.
\end{enumerate}

في هذه العملية، ملف المواصفات هو «مصدر الحقيقة الوحيد» — يمكن إعادة توليد الكود، لكن المواصفات يجب أن تكون صحيحة أولًا.

كتابة ملف المواصفات

هل المواصفات مكتملة؟

مكتملة

الذكاء الاصطناعي يولّد الكود وفق المواصفات

تشغيل الاختبارات للتحقق من الناتج

غير مكتملة / اختبارات لا تنجح

مشكلة في الكود صحح الكود؛ مشكلة في المواصفات صحح المواصفات

العودة إلى التحقق

كل الاختبارات ناجحة تسليم

\begin{center}\small\color{gray}الشكل 3: سير عمل التطوير المدفوع بالمواصفات (SDD)\end{center}

تكمن ميزة SDD في أن ملف المواصفات أكثر استقرارًا من الشيفرة. فعندما تتغير النماذج الأساسية أو الأدوات، يبقى ملف المواصفات لديك صالحًا، ويكفي أن تطلب من أداة الذكاء الاصطناعي الجديدة إعادة توليد الشيفرة. ولهذا يؤكد هذا الكتاب على أن «المنهجية أهم من الأدوات» — فملف المواصفات هو التجسيد الملموس لمنهجيتك.

\booksection{ثالثًا: إدارة نافذة السياق}

حتى مع ملف ذاكرة المشروع، ستواجه مشكلة «نقص السياق» في المحادثات. نافذة سياق النموذج (كمية المعلومات التي يعالجها في وقت واحد) محدودة. عندما تطول المحادثة وتكثر الملفات المذكورة، «ينسى» الذكاء الاصطناعي المعلومات المبكرة، فتنخفض جودة المخرجات — كما لو دردشت مع صديق ست ساعات، فنسي أول ما قلته.

الاستراتيجية الأساسية لإدارة نافذة السياق هي «الاستشهاد المُصعَّد»:

\begin{center}
\begin{tabularx}{\linewidth}{|X|X|X|}
\hline
\textbf{طبقة السياق} & \textbf{المحتوى} & \textbf{طريقة الإدارة} \tabularnewline
\hline
الطبقة الأولى: ذاكرة المشروع & المكدس التقني، والمواصفات، وبنية المجلدات & تُكتب في ملفات مثل CLAUDE.md، وتُحمَّل تلقائيًا دون تكرار \tabularnewline
الطبقة الثانية: استشهاد الملفات & ملفات الكود المتعلقة بالمهمة الحالية & استشهد بدقة برمز @، ولا تلصق الملف كاملًا في المحادثة \tabularnewline
الطبقة الثالثة: سياق المحادثة & سجل محادثة المهمة الحالية & بعد إتمام مهمة، افتح محادثة جديدة لتجنب تضخم السياق \tabularnewline
الطبقة الرابعة: الأمر الفوري & متطلبات الخطوة الحالية & واضح ومختصر ومحدد، وتجنب الوصف المطول \tabularnewline
\hline
\end{tabularx}
\end{center}

\begin{center}\small\color{gray}استراتيجية الإدارة المُصعَّدة لنافذة السياق\end{center}

\booksection{رابعًا: نقاط عملية في هندسة السياق}

هندسة السياق ليست إعدادًا يُجرى مرة واحدة، بل عملية تحسين مستمرة. وإليك بعض النقاط التي ثبتت فعاليتها مرارًا وتكرارًا في الممارسة:

\begin{enumerate}
  \item حدِّث ملف الذاكرة بالتزامن مع تطور المشروع. فعند إدخال تقنية جديدة أو تغيير البنية، حدِّث ملف الذاكرة فورًا بدلًا من الشرح المؤقت داخل المحادثة.
  \item أحسن استخدام الإحالة بـ @ بدلًا من لصق النص كاملًا. تدعم معظم أدوات البرمجة بالذكاء الاصطناعي صيغة @file أو @folder التي تُدخل محتوى الملف إلى سياق الذكاء الاصطناعي في الخلفية، وهي أكثر كفاءة من اللصق اليدوي ولا تُلوِّث المحادثة.
  \item اجعل المحادثات «قصيرة ومركزة». كل محادثة تعالج مهمة واحدة أو نقطة وظيفية واحدة، وبعد الانتهاء افتح محادثة جديدة للمهمة التالية. فالمحادثات الطويلة لا تستهلك نافذة السياق فحسب، بل تُشتت تركيز الذكاء الاصطناعي أيضًا.
  \item دَع الذكاء الاصطناعي «يرى» النتائج بدلًا من «وصفها» له. فعند مواجهة خطأ (bug)، أشر إلى سجل الأخطاء مباشرة بـ @ بدلًا من وصفه بكلماتك مثل «يبدو أن هناك خطأً ما». فرؤية الذكاء الاصطناعي للمعلومة الأصلية أنجع من رؤيته لروايتك عنها.
  \item نظِّف والأرشف دوريًا. فالمحادثات المنتهية والملفات المؤقتة تشغل مساحة العمل. نظِّف المحتوى غير ذي الصلة بانتظام وحافظ على نظافة مجلد المشروع، فهذا يساعد الذكاء الاصطناعي أيضًا على فهم بنية المشروع بكفاءة أكبر.
\end{enumerate}

\begin{tcbpractice}{هيّا نطبق: اكتب لمشروعك ملف ذاكرة من 70 سطرًا}
اختر مشروعًا تعمل عليه الآن (أو تنوي العمل عليه)، واكتب ملف CLAUDE.md يتضمن: نظرة عامة على المشروع (من جملتين إلى ثلاث)، والحزمة التقنية (اذكرها)، ومعايير الشيفرة (من 3 إلى 5 بنود)، وهيكل المجلدات (ارسمه شجرة)، واتفاقيات التطوير (من بندين إلى ثلاثة). ويُفضل ألا يتجاوز 70 سطرًا — وستجد بعد الانتهاء أن فهمك للمشروع صار أوضح بكثير.

\end{tcbpractice}

\begin{tcbsummary}{خلاصة الفصل}
\begin{itemize}
  \item ذاكرة الذكاء الاصطناعي «قصيرة المدى»: في كل محادثة جديدة تبدأ معرفته بالمشروع من الصفر
  \item ملف ذاكرة المشروع (مثل CLAUDE.md) هو «الذاكرة طويلة المدى» للذكاء الاصطناعي، ويُحمَّل تلقائيًا عند كل محادثة
  \item ملف الذاكرة الجيد = نظرة عامة + حزمة تقنية + معايير + هيكل المجلدات + اتفاقيات، و70 سطرًا تكفي
  \item التطوير الموجَّه بالمواصفات (SDD): اكتب المواصفات قبل الشيفرة، فالمواصفات هي «المصدر الوحيد للحقيقة» وهي أكثر استقرارًا من الشيفرة
  \item يُدار السياق في أربع طبقات: ذاكرة المشروع إحالة الملفات بـ @ سجل المحادثة التعليمات الفورية
  \item خمس نقاط عملية: التحديث المتزامن، والإحالة بـ @، والمحادثات القصيرة، وإعطاء المعلومة الأصلية، والتنظيف الدوري
\end{itemize}

\end{tcbsummary}
